\section{问题4：数据采集建议与策略闭环}

基于前三问的建模分析发现，建议游戏服务器优先采集以下两类新埋点数据。

\subsection{建议一：社交行为数据}

\textbf{建议采集字段：}好友互动频次（好友列表、私聊次数、赠礼次数）、联盟聊天消息数及活跃时段、组队（集结/集结参与）次数与成功率、被攻击/被侦查记录。

\textbf{理由：}问题1的Cox回归显示``is\_in\_league''是留存的最强保护因子之一（$\exp(\beta) = 0.90$, $p < 0.001$），问题1的分段KM分析进一步表明前3天加联盟用户的7日留存率（28.0\%）是未加联盟用户（1.2\%）的23倍。然而，当前仅有布尔型``是否加联盟''标签，无法量化社交深度。精确的社交行为数据可将``浅社交''（仅加入联盟但不活跃）与``深社交''（频繁互动、组队、聊天）用户区分开来，二者的留存和付费行为可能存在数量级差异。

\textbf{策略闭环修正：}
\begin{itemize}
    \item 社交深度高 $+$ 进度压力低 → 低频推送（$\leq$1次/周），减少打扰
    \item 社交深度高 $+$ 进度压力高 → 中频推送（2次/周），推送联盟专属礼包
    \item 社交深度低 $+$ 进度压力低 → 中频推送，引导加入联盟
    \item 社交深度低 $+$ 进度压力高 → 精准触发式推送，低价社交激励礼包
\end{itemize}

\subsection{建议二：实时战力/进度数据}

\textbf{建议采集字段：}PVP/GVG战斗胜负记录及战力变化、各建筑等级及升级剩余时间、科技树研究进度、主线/支线任务完成进度、英雄收集数量与星级分布。

\textbf{理由：}问题2的资源-等级弹性分析显示资源消耗与等级增长相关性显著（$r = 0.39$--$0.50$），但当前数据仅有资源变化记录，缺少``进度目标''信息。玩家是否有未完成的升级/研究/任务是判断资源缺口的关键——资源缺口率应定义为``（所需资源 - 持有资源）/ 所需资源''，仅凭消耗量无法精确计算。PVP胜负记录则是判断``进度焦虑''的核心信号，连续N次PVP失败是触发``战力突破礼包''的最佳时机。

\textbf{策略闭环修正：}
\begin{itemize}
    \item 建立``进度压力指数''：综合PVP胜率、建筑排队时间、资源缺口率三维指标
    \item 资源缺口率 $> 50\%$ → 推送低价补给包（6元，降低决策门槛）
    \item 资源缺口率 $> 80\%$ 且 PVP连败 $\geq 3$次 → 推送高价突破包（68--128元）
    \item 升级受阻（建筑等待时间 $>$ 阈值）→ 推送加速道具包
\end{itemize}

\subsection{AB测试验证框架}

新服上线后，建议采用AB测试验证增强策略的有效性：
\begin{itemize}
    \item \textbf{对照组：}仅使用问题3的基准方案（固定礼包策略）
    \item \textbf{实验组：}使用社交$+$进度维度的增强动态策略
    \item \textbf{评估指标：}对比两组30日ARPU、留存率、各聚类付费转化率
    \item \textbf{迭代机制：}每周根据数据更新策略参数，持续优化
\end{itemize}
